OpenFisca-France 176.1 (législation des revenus 2024) ; DGFiP, IRCOM revenus 2024 ; calculs de l'auteur.\par Note : exercice 3 : taux marginal de 13,0 \% au premier euro, 16,0 \% à 20 000 \euro{}, 19,2 \% à 80 000 \euro{}, 50 \% à 600 000 \euro{} par adulte (revenu après cotisations), en ligne droite entre ces seuils. Impôt individuel sur le revenu après cotisations, franchise de 6 000 \euro{} par enfant déduite du revenu du parent qui gagne le plus ; même recette qu'aujourd'hui. Salariés, IR + CSG + CRDS en \% du salaire brut ; taux marginal sur 1 000 \euro{} de plus, lissé sur 2 000 \euro{} (médiane glissante) pour effacer les pointes isolées. Aujourd'hui, le taux marginal du couple monte à 21 \% vers 60 000 \euro{} puis retombe à 18 \% entre 86 000 et 95 000 \euro{} : en dessous, la sortie progressive de la décote alourdit l'IR.
